\section{Разложение булевых функций по переменным. Нормальные формы. КНФ, ДНФ, СКНФ, СДНФ. Сокращенные дизъюнктивные формы.}

\subsection*{Разложение по переменным (Шеннон)}
Любая функция может быть представлена через свои значения на частных наборах:
$$f(x_1, \dots, x_n) = (x_i \land f(x_1, \dots, 1, \dots, x_n)) \lor (\neg x_i \land f(x_1, \dots, 0, \dots, x_n)).$$
Если применить это разложение последовательно по всем $n$ переменным, мы придем к понятию совершенных нормальных форм.

\subsection*{Совершенные нормальные формы}

\paragraph{Определение СДНФ}
\textbf{Совершенная дизъюнктивная нормальная форма} — это реализация функции в виде дизъюнкции элементарных конъюнкций, каждая из которых содержит \textbf{все} переменные (в прямом или инверсном виде).
$$f(x_1, \dots, x_n) = \bigvee_{\sigma: f(\sigma)=1} (x_1^{\sigma_1} \land \dots \land x_n^{\sigma_n}).$$
\textit{Процедура:} Берем все строки таблицы истинности, где $f=1$. Для каждой строки пишем конъюнкцию: если в строке $x_i=1$, пишем $x_i$; если $x_i=0$, пишем $\neg x_i$.

\paragraph{Определение СКНФ}
\textbf{Совершенная конъюнктивная нормальная форма} — это реализация функции в виде конъюнкции элементарных дизъюнкций, содержащих все переменные.
$$f(x_1, \dots, x_n) = \bigwedge_{\sigma: f(\sigma)=0} (x_1^{\neg\sigma_1} \lor \dots \lor x_n^{\neg\sigma_n}).$$
\textit{Процедура:} Берем строки, где $f=0$. Пишем дизъюнкцию: если в строке $x_i=0$, пишем $x_i$; если $x_i=1$, пишем $\neg x_i$.

\paragraph{Теорема о единственности}
Всякая булева функция (кроме констант 0 и 1 в соответствующих случаях) имеет **единственную** СДНФ и СКНФ (при установленном лексикографическом порядке переменных и слагаемых).

\subsection*{Эквивалентные преобразования к СДНФ}
Любую формулу можно привести к СДНФ за 7 шагов:
1. \textbf{Элиминация:} замена всех операций ($\to, \oplus, \equiv$) на $\{\lor, \land, \neg\}$.
2. \textbf{Протаскивание отрицаний:} по законам де Моргана $\neg$ доводится до переменных.
3. \textbf{Раскскрытие скобок:} использование дистрибутивности конъюнкции.
4. \textbf{Удаление нулей:} убирание конъюнкций вида $(x \land \neg x)$.
5. \textbf{Приведение подобных:} удаление лишних вхождений по идемпотентности.
6. \textbf{Расщепление:} добавление недостающих переменных по правилу $A = (A \land x) \lor (A \land \neg x)$.
7. \textbf{Сортировка:} упорядочивание переменных внутри конъюнкций и самих конъюнкций.

\subsection*{ДНФ и Сокращенная форма}

\paragraph{ДНФ (Дизъюнктивная нормальная форма)}
Это дизъюнкция элементарных конъюнкций произвольного ранга. В отличие от СДНФ, здесь в слагаемых могут отсутствовать некоторые переменные.

\paragraph{Сокращенная ДНФ}
Это ДНФ, состоящая из суммы всех \textbf{простых импликант} функции.
\begin{itemize}
    \item \textbf{Импликанта:} конъюнкция $K$, такая что $K \to f = 1$.
    \item \textbf{Простая импликанта:} импликанта, из которой нельзя удалить ни одной переменной, чтобы она осталась импликантой.
\end{itemize}
Сокращенная ДНФ получается из СДНФ путем всех возможных операций **склеивания** ($(A \land x) \lor (A \land \neg x) = A$) и последующего поглощения.
